%==============================================================
%  EV Final Poster -- A0 portrait  (beamerposter)
%  Dual-pillar: identifiability = observation x excitation
%  Layout / palette / fonts after EV海報參考模板.pptx
%  Build:  python3 refresh.py && latexmk -pdf main.tex
%==============================================================
\documentclass[final,t]{beamer}
\usepackage[size=a0,orientation=portrait,scale=1.30]{beamerposter}

\usepackage[utf8]{inputenc}
\usepackage[T1]{fontenc}
\usepackage{helvet}            % Arial-like sans
\renewcommand{\familydefault}{\sfdefault}
\usepackage{amsmath,amssymb}
\usepackage{graphicx}
\usepackage{booktabs}
\usepackage{tikz}
\usetikzlibrary{arrows.meta,positioning,fit}

\input{generated/results.tex}

%---------------- palette --------------------------------------
\definecolor{pblue}{HTML}{2E5394}
\definecolor{ptint}{HTML}{EEF3F9}
\definecolor{good}{HTML}{1B7F3B}
\definecolor{bad}{HTML}{B11A1A}

\setbeamercolor{block title}{bg=pblue,fg=white}
\setbeamercolor{block body}{bg=ptint,fg=black!88}
\setbeamerfont{block title}{series=\bfseries,size=\LARGE}
\setbeamercolor{itemize item}{fg=pblue}
\setbeamercolor{itemize subitem}{fg=pblue}
\setbeamertemplate{itemize items}[circle]

\setlength{\parskip}{0.24\baselineskip}
\usepackage{etoolbox}
\AtBeginEnvironment{itemize}{\setlength{\itemsep}{0.6ex}\setlength{\parskip}{0pt}}
\AtBeginEnvironment{enumerate}{\setlength{\itemsep}{0.6ex}\setlength{\parskip}{0pt}}

% generous block-title padding, compact gaps between blocks
\setbeamertemplate{block begin}{
  \vskip0.4ex
  \begin{beamercolorbox}[colsep*=0pt,sep=1.8ex,rounded=true]{block title}
    \hspace{0.5cm}\usebeamerfont{block title}\insertblocktitle
  \end{beamercolorbox}
  \nointerlineskip\vskip0.6ex
  \usebeamerfont{block body}%
  \begin{beamercolorbox}[colsep*=0pt,sep=1.4ex,rounded=true,vmode]{block body}
    \vskip0.3ex
}
\setbeamertemplate{block end}{
    \vskip0.3ex
  \end{beamercolorbox}
}

\newcommand{\cmark}{\textcolor{good}{\textbf{\checkmark}}}
\newcommand{\xmark}{\textcolor{bad}{$\boldsymbol{\times}$}}
\newcommand{\oracle}{\textcolor{bad}{\footnotesize\,(oracle)}}
\newcommand{\sandbox}{\textcolor{bad}{\footnotesize\,(forward+FD, synthetic)}}
\newcommand{\imgloss}{\textcolor{pblue}{\textbf{image loss}}}

\setlength{\columnsep}{2cm}

\begin{document}
\begin{frame}[t]{}

%================= HEADER ======================================
\begin{beamercolorbox}[wd=\paperwidth,sep=1.0cm,leftskip=2cm,rightskip=2cm]{block title}
  \centering
  {\Huge\bfseries Probing Physical-Parameter Identifiability from Motion}\\[3mm]
  {\LARGE Observation $\times$ Excitation, via Differentiable MPM \texttt{+} 3D Gaussian Splatting}\\[3mm]
  {\Large HONG, Casey (b10401006)\ \ $\cdot$\ \ Embodied Vision (CSIE5421), NTU}
\end{beamercolorbox}
\vspace{0.15cm}

%================= ROW 1 : Question | Takeaways ===============
\begin{columns}[t,totalwidth=\textwidth]
\begin{column}{0.49\textwidth}
\begin{block}{The question \& three regimes}
\textbf{When is a physical parameter recoverable from how an object moves?}
We probe this along \textbf{two orthogonal axes} over the differentiable map
$g=R_{\text{3DGS}}\!\circ\Phi_{\text{MPM}}$:

\vspace{2mm}\centering
\begin{tikzpicture}[font=\large,node distance=10mm,
  box/.style={draw=pblue,line width=1.2pt,rounded corners,fill=white,
              minimum height=1.8cm,minimum width=2.3cm,inner sep=4pt,align=center},
  >={Stealth[length=6mm]}]
  \node[box] (y) {$Y$\\[0.5mm]\small $E,v_0,F_0$};
  \node[box,right=of y] (mpm) {$\Phi_{\text{MPM}}$};
  \node[box,right=of mpm] (gs) {$R_{\text{3DGS}}$};
  \node[box,right=of gs] (x) {$X$\\[0.5mm]\small video};
  \draw[->,line width=1.2pt] (y)--(mpm);
  \draw[->,line width=1.2pt] (mpm)--(gs);
  \draw[->,line width=1.2pt] (gs)--(x);
\end{tikzpicture}

\vspace{2mm}\raggedright
Three supervision regimes, kept honestly separate:
\begin{itemize}
  \item \imgloss{} on pixels --- \textbf{recoverable from video} (Pillar A).
  \item \textcolor{bad}{trajectory loss} --- needs GT motion, \emph{oracle} only.
  \item \textcolor{bad}{forward $+$ finite difference} --- synthetic identifiability
        analysis (Pillar B).
\end{itemize}
\end{block}
\end{column}

\begin{column}{0.49\textwidth}
\begin{block}{Key Takeaways}
\begin{enumerate}
  \item Identifiability from motion $=$ \textbf{observation geometry}
        $\times$ \textbf{excitation design}.
  \item \textbf{(A) Observation:} from \imgloss, single-view $(E,v_0)$ is
        degenerate; a \textbf{second view resolves both} to $\sim1\%$.
  \item \textbf{(B) Excitation:} $E$ enters motion via a \emph{frequency}
        channel (displacement-excited $\Rightarrow$ spectral loss) or an
        \emph{amplitude} channel (force-excited $\Rightarrow$ time loss) ---
        explaining why gravity/contact scenes are easy to fit.
  \item The pre-stress $F_0$ is \textbf{itself identifiable} (global 6-DOF,
        coarse 15-DOF field). \textbf{Future:} a generative physics prior (VSD).
\end{enumerate}
\end{block}
\end{column}
\end{columns}

%================= ROW 2 : PILLAR A (full) ====================
\begin{block}{Pillar A --- Observation geometry \ \textnormal{\normalsize(image loss, recoverable from video)}}
\begin{columns}[T,totalwidth=\textwidth]
%--- A.1 recovery loop
\begin{column}{0.30\textwidth}
\centering
\textbf{Recovery loop.} Freeze the 3DGS scene, leave only $Y$ free, backprop
the image loss through render \texttt{+} a differentiable MPM solver.

\vspace{4mm}
\begin{tikzpicture}[font=\large,>={Stealth[length=6mm]},
  st/.style={draw=pblue,line width=1.2pt,rounded corners,minimum height=2.2cm,
             inner sep=6pt,fill=white,align=center}]
  \node[st] (init) {init $Y$};
  \node[st,right=10mm of init] (fwd) {$\Phi\!\to\!R$};
  \node[st,right=10mm of fwd] (loss) {$\|\hat X\!-\!X\|$};
  \draw[->,line width=1.2pt] (init)--(fwd);
  \draw[->,line width=1.2pt] (fwd)--(loss);
  \draw[->,line width=1.5pt,pblue] (loss.north) to[bend right=34]
        node[above,font=\small,pblue]{update $Y$} (init.north);
\end{tikzpicture}

\vspace{4mm}\raggedright
\textbf{Two levers:} \emph{observation} (\#views) and \emph{excitation design}.
$Y$ scales from a global scalar $E$ to per-particle fields.
\end{column}

%--- A.2 the identifiability table
\begin{column}{0.36\textwidth}
\centering
\textbf{A second view breaks the $E$--$v_0$ degeneracy}\\[3mm]
\renewcommand{\arraystretch}{1.45}\setlength{\tabcolsep}{10pt}
{\large
\begin{tabular}{lcc}
\toprule
 & \textbf{known $v_0$} & \textbf{unknown $(E,v_0)$}\\
\midrule
\textbf{single view} & $E$ \fitRgbErr\% \,\cmark
   & $E{\to}\svBestE$ ($-\svErrE\%$) \,\xmark\\
\textbf{multi view}  & ---
   & $E$ \mvErrE\%, $v_0$ \mvErrVz\% \,\cmark\\
\bottomrule
\end{tabular}}

\vspace{4mm}\raggedright\normalsize
One view cannot separate ``stiffer\,$+$\,faster'' from ``softer\,$+$\,slower''
(they render alike), so the joint fit \textbf{collapses}; a side view supplies
the missing constraint ($\sim\mvWallMin$\,min). \textbf{Excitation design}
likewise gates a recoverable $v_0$ field: a \emph{bend} gives rel-$L_2$
\fieldBendRelL{} vs a harder \emph{mid-kick} \fieldMidRelL.
\end{column}

%--- A.3 figures
\begin{column}{0.30\textwidth}
\centering
\includegraphics[width=0.92\linewidth]{figures/fit_rgb_panel.png}\\[1mm]
{\small A. known $v_0\Rightarrow E$ \fitRgbErr\%}\\[2.5mm]
\includegraphics[width=0.475\linewidth]{figures/sv_fail_panel.png}\hfill
\includegraphics[width=0.475\linewidth]{figures/field_bend_quiver.png}\\[1mm]
{\small B./C. single fail vs multi solve\ \ $\cdot$\ \ D. $v_0$ bend field}
\end{column}
\end{columns}
\end{block}

%================= ROW 3 : PILLAR B (full) ====================
\begin{block}{Pillar B --- Excitation design \& the two $E$ channels \ \textnormal{\normalsize\textcolor{bad}{(forward $+$ finite-difference analysis; synthetic --- autograd/image-loss recovery in progress)}}}
\begin{columns}[T,totalwidth=\textwidth]
%--- B.1 channel table
\begin{column}{0.29\textwidth}
\centering
\renewcommand{\arraystretch}{1.5}\setlength{\tabcolsep}{6pt}
{\normalsize
\begin{tabular}{lll}
\toprule
\textbf{excitation} & \textbf{$E$ channel} & \textbf{loss}\\
\midrule
displacement & frequency & spectral\\
(release)    & $\propto\!\sqrt{E}$ & \\[1mm]
force        & amplitude & time\\
(drop/contact) & $\propto\!1/E$ & \\
\bottomrule
\end{tabular}}

\vspace{4mm}\raggedright\normalsize
The excitation \emph{type} decides which channel carries $E$, hence which loss
can read it --- \textbf{no universal loss}. This explains why gravity/contact
benchmarks fit with a plain time/image loss, while pre-stress release needs a
\emph{spectral} loss (time-domain saturates into a beat plateau).
\end{column}

%--- B.2 channel landscapes + signals
\begin{column}{0.40\textwidth}
\centering
\includegraphics[width=0.485\linewidth]{figures/f0_chan_release.png}\hfill
\includegraphics[width=0.485\linewidth]{figures/f0_chan_drop.png}\\[1mm]
{\small Complementary landscapes: time flat on \emph{release}, spectral flat on
\emph{drop}.}\\[3mm]
\includegraphics[width=0.86\linewidth]{figures/f0_signals.png}\\[1mm]
{\small Displacement-excited: time-domain trajectories saturate; the spectrum
stays discriminative.}
\end{column}

%--- B.3 F0 recovery + Nyquist
\begin{column}{0.29\textwidth}
\centering
\textbf{$F_0$ is itself identifiable}\\[2mm]
\includegraphics[width=0.92\linewidth]{figures/f0_S_recovery.png}\\[1mm]
{\small Global 6-DOF $V_0=\exp(S)$: all components recovered, max$|$err$|$
\sErrMax. Coarse 15-DOF field also recovered (one shallow bend degeneracy mapped).}\\[3mm]
\includegraphics[width=0.78\linewidth]{figures/f0_period.png}\\[1mm]
{\small Caveat: period $\propto1/\sqrt{E}$ hits a Nyquist wall $\Rightarrow$
an aliasing false basin at high $E$.}
\end{column}
\end{columns}
\end{block}

%================= ROW 4 : Extensions + Future (full) =========
\begin{block}{Extensions \& toward generative physics}
\begin{columns}[c,totalwidth=\textwidth]
\begin{column}{0.66\textwidth}
\begin{itemize}
  \item \textbf{Amortised multi-impulse:} one shared $E$ from several excitations
        --- \textbf{\multiobsErr\%} (vs.\ \multiobsSingleErr\% single). Shared
        material is learnable. \oracle
  \item \textbf{Non-uniform $E$ field} and \textbf{$F_0$-conditioned} recovery
        ($E$ on a pre-stressed rest state) --- newest direction. \oracle\sandbox
  \item \textbf{Future:} amortise a \emph{scene-conditioned physics prior} and
        sample plausible physics via VSD $+$ reconstruction guidance.
\end{itemize}
\end{column}
\begin{column}{0.32\textwidth}
\centering
\includegraphics[width=0.62\linewidth]{figures/ext_multiobs.png}\\[1mm]
{\small Shared-$E$ convergence across impulses.}
\end{column}
\end{columns}
\end{block}

\end{frame}
\end{document}
